\chapter{Dalsze własności pochodnych}
\label{ma1:ch:derivative-properties}

Pochodna nie musi być ciągła (przykład \ref{ma1:ex:discontinuous-derivative}), lecz ma własność wartości pośredniej. Twierdzenia o wartości średniej z rozdziału \ref{ma1:ch:mean-value-theorems} pozwalają także wyznaczać granice pewnych ilorazów funkcji.

\section{Twierdzenie Darboux}
\begin{twierdzenie}[Darboux]
\label{ma1:thm:derivative-darboux}
Niech $f:I\to\R$ będzie różniczkowalna na otwartym przedziale $I$. Jeżeli $a,b\in I$, $a<b$, oraz $\lambda$ leży ściśle między $f'(a)$ i $f'(b)$, to istnieje $c\in(a,b)$ takie, że $f'(c)=\lambda$. W szczególności obraz każdego podprzedziału $I$ przez $f'$ jest przedziałem.
\end{twierdzenie}
\begin{proof}
Załóżmy najpierw, że $f'(a)<\lambda<f'(b)$. Niech $H(x)=f(x)-\lambda x$. Funkcja $H$ jest różniczkowalna na $I$, zatem ciągła na $[a,b]$. Jej pochodna spełnia
\[
 H'(a)=f'(a)-\lambda<0,\qquad H'(b)=f'(b)-\lambda>0.
\]
Z definicji pochodnej dla dostatecznie małego $h>0$, $h<b-a$, mamy $(H(a+h)-H(a))/h<0$, czyli $H(a+h)<H(a)$. Podobnie $(H(b-h)-H(b))/(-h)>0$, skąd $H(b-h)<H(b)$. Żaden koniec nie jest więc punktem osiągania minimum na $[a,b]$.

Z twierdzenia o osiąganiu kresów \ref{ma1:cor:interval-extrema} istnieje punkt $c\in[a,b]$, w którym $H$ osiąga minimum. Musi on należeć do $(a,b)$. Jest zatem punktem minimum lokalnego, a twierdzenie Fermata \ref{ma1:thm:fermat} daje $H'(c)=0$, czyli $f'(c)=\lambda$.

Jeżeli $f'(a)>\lambda>f'(b)$, stosujemy dowiedziony przypadek do $-f$ i liczby $-\lambda$. Otrzymujemy $(-f)'(c)=-\lambda$, czyli $f'(c)=\lambda$.

Niech teraz $J\subseteq I$ będzie niepustym przedziałem i niech $u,v\in f'(J)$, $u<v$. Wybieramy punkty $x,y\in J$ takie, że $f'(x)=u$, $f'(y)=v$. Ponieważ $x\ne y$, odcinek o końcach $x,y$ leży w $J$. Każda liczba z $(u,v)$ jest wartością pochodnej w jego wnętrzu na mocy pierwszej części twierdzenia. Zatem $[u,v]\subseteq f'(J)$, co dowodzi, że $f'(J)$ jest przedziałem.
\end{proof}

\begin{center}
\begin{tikzpicture}[x=2.1cm,y=1.3cm,>=Stealth]
 \draw[->] (-.25,0)--(3.35,0) node[right] {$x$};
 \draw[->] (0,-.2)--(0,2.9) node[above] {$y$};
 \draw[blue!65!black,thick,domain=.25:2.75,samples=100] plot (\x,{(\x-1.5)*(\x-1.5)+.5});
 \draw[orange!85!black,thick] (.95,.5)--(2.05,.5);
 \draw[gray,dotted] (.5,0)--(.5,1.5);
 \draw[gray,dotted] (1.5,0)--(1.5,.5);
 \draw[gray,dotted] (2.5,0)--(2.5,1.5);
 \fill (.5,1.5) circle (1.6pt);
 \fill (1.5,.5) circle (1.6pt);
 \fill (2.5,1.5) circle (1.6pt);
 \node[below] at (.5,0) {$a$};
 \node[below] at (1.5,0) {$c$};
 \node[below] at (2.5,0) {$b$};
 \node[blue!65!black,above left] at (.65,2.15) {$H$};
 \node[above] at (1.5,.5) {$H'(c)=0$};
 \node[above right] at (.5,1.5) {$H'(a)<0$};
 \node[above left] at (2.5,1.5) {$H'(b)>0$};
\end{tikzpicture}

\small Rysunek: po odjęciu funkcji liniowej $\lambda x$ szukamy punktu o poziomej stycznej. Przeciwne znaki pochodnej na końcach wymuszają osiągnięcie minimum we wnętrzu odcinka.
\end{center}

\begin{wniosek}
\label{ma1:cor:derivative-no-jumps}
Jeżeli funkcja $f$ jest różniczkowalna na otwartym przedziale $I$ i $a\in I$, to każda istniejąca jednostronna granica $f'$ w $a$ (również rozpatrywana w rozszerzonej prostej rzeczywistej) jest skończona i równa $f'(a)$.
\end{wniosek}
\begin{proof}
Wybieramy $b>a$ w $I$. Z różniczkowalności $f$ w $a$ wynika jej ciągłość w $a$ oraz istnienie skończonej pochodnej prawostronnej równej $f'(a)$. Twierdzenie \ref{ma1:thm:derivative-endpoint-limit} zastosowane na $[a,b]$ mówi, że jeśli granica prawostronna $f'$ istnieje, to jest równa tej pochodnej, a więc nie może być nieskończona. Dla granicy lewostronnej stosujemy wersję twierdzenia dla prawego końca odcinka $[d,a]\subset I$, gdzie $d<a$.
\end{proof}
W szczególności pochodna nie może mieć nieciągłości skokowej ani usuwalnej w punkcie wewnętrznym. Nie oznacza to ciągłości: przykład \ref{ma1:ex:discontinuous-derivative} pokazuje, że granica pochodnej może nie istnieć. Własność Darboux daje dodatkowo informację o przyjmowaniu wszystkich wartości pomiędzy dwiema wartościami pochodnej, niezależnie od zachowania granic jednostronnych.

\begin{przyklad}
\label{ma1:ex:darboux-step}
Funkcja $g:\R\to\R$ określona przez $g(x)=-1$ dla $x<0$ oraz $g(x)=1$ dla $x\ge0$ nie jest pochodną żadnej funkcji różniczkowalnej na $\R$. Gdyby $g=f'$, to z $f'(-1)=-1$ i $f'(1)=1$ twierdzenie Darboux dla $\lambda=0$ dawałoby punkt $c\in(-1,1)$ z $g(c)=0$, którego nie ma.
\end{przyklad}

\section{Reguła de l'Hospitala dla postaci $0/0$}
\begin{twierdzenie}[Reguła de l'Hospitala, postać $0/0$]
\label{ma1:thm:lhopital-zero}
Niech $I$ będzie otwartym przedziałem, $a$ jednym z jego końców, skończonym lub nieskończonym, a $f,g:I\to\R$ funkcjami różniczkowalnymi. Zakładamy, że $g'(x)\ne0$ w pewnym jednostronnym otoczeniu końca $a$ w $I$ oraz
\[
 \lim_{x\to a}f(x)=\lim_{x\to a}g(x)=0,\qquad
 \lim_{x\to a}\frac{f'(x)}{g'(x)}=L
 \quad (L\in\R\cup\{-\infty,+\infty\}).
\]
Wtedy $g(x)\ne0$ dostatecznie blisko $a$ oraz $\lim_{x\to a} f(x)/g(x)=L$.
\end{twierdzenie}
\begin{proof}
Ograniczamy $I$ do przedziału przylegającego do końca $a$, na którym $g'\ne0$. Z twierdzenia Rolle'a \ref{ma1:thm:rolle} wynika, że dla różnych $x,y$ z tego przedziału mamy $g(x)\ne g(y)$; inaczej istniałoby zero $g'$ pomiędzy nimi. W szczególności $g$ ma co najwyżej jedno zero. Jeżeli ono istnieje, możemy jeszcze zmniejszyć otoczenie końca $a$, aby je wykluczyć.

Dla $x$ w tym otoczeniu i $y$ leżącego ściśle między $x$ i końcem $a$ wniosek \ref{ma1:cor:cauchy-ratio} daje punkt $c$ między $x$ i $y$ taki, że
\[
 \frac{f(x)-f(y)}{g(x)-g(y)}=\frac{f'(c)}{g'(c)}.
\]
Założenia wniosku są spełnione: funkcje są ciągłe na odcinku o końcach $x,y$, różniczkowalne w jego wnętrzu i $g'$ jest tam niezerowa.

Jeżeli $L$ jest skończona, to dla $\varepsilon>0$ wybieramy jeszcze mniejsze otoczenie końca $a$, w którym $|f'/g'-L|<\varepsilon/2$. Dla każdego $x$ z tego otoczenia oraz każdego $y$ pomiędzy $x$ i $a$ również punkt $c$ do niego należy. Stąd
\[
 L-\varepsilon/2<
 \frac{f(x)-f(y)}{g(x)-g(y)}
 <L+\varepsilon/2.
\]
Przy ustalonym $x$ przechodzimy do granicy $y\to a$. Z założeń $f(y),g(y)\to0$ i $g(x)\ne0$ wynika, że środkowy iloraz dąży do $f(x)/g(x)$. Otrzymujemy $|f(x)/g(x)-L|\le\varepsilon/2<\varepsilon$ dla wszystkich $x$ w wybranym otoczeniu. To dowodzi tezy.

Jeżeli $L=+\infty$, dla dowolnego $K>0$ wybieramy otoczenie, w którym $f'/g'>K+1$. Ten sam iloraz różnic jest większy od $K+1$, a po przejściu do granicy $y\to a$ dostajemy $f(x)/g(x)\ge K+1>K$. Dla $L=-\infty$ wybieramy otoczenie z $f'/g'<-K-1$ i otrzymujemy $f(x)/g(x)\le-K-1<-K$. Argument nie zależy od tego, czy koniec $a$ jest lewy, prawy, skończony czy nieskończony.
\end{proof}
Dla granicy dwustronnej w punkcie skończonym stosujemy twierdzenie oddzielnie z lewej i z prawej strony. Jeżeli obie granice są równe tej samej liczbie $L$ (lub tej samej nieskończoności), istnieje granica dwustronna i wynosi $L$. Nie jest wymagana różniczkowalność funkcji w samym punkcie granicznym.

\section{Reguła de l'Hospitala przy nieskończonym mianowniku}
\begin{twierdzenie}[Reguła de l'Hospitala]
\label{ma1:thm:lhopital-infinity}
Przyjmijmy takie same założenia o $I,a,f,g$ oraz o niezerowości $g'$ jak w twierdzeniu \ref{ma1:thm:lhopital-zero}. Jeżeli $|g(x)|\to+\infty$ przy $x\to a$ oraz
\[
 \lim_{x\to a}\frac{f'(x)}{g'(x)}=L
 \quad (L\in\R\cup\{-\infty,+\infty\}),
\]
to $\lim_{x\to a}f(x)/g(x)=L$.
\end{twierdzenie}
\begin{proof}
Nie musimy zakładać, że licznik ma granicę nieskończoną. Mianownik $g$ jest ostatecznie niezerowy z założenia o jego module.

Jeżeli $L$ jest skończona, dla $\varepsilon>0$ wybieramy przedział przylegający do $a$, na którym $g'\ne0$ i $|f'/g'-L|<\varepsilon/4$. Ustalamy w nim punkt $b$. Dla każdego $x$ pomiędzy $b$ i końcem $a$, wniosek \ref{ma1:cor:cauchy-ratio} daje
\[
 Q(x):=\frac{f(x)-f(b)}{g(x)-g(b)}
 =\frac{f'(c)}{g'(c)}
\]
dla pewnego $c$ między $b$ i $x$. W szczególności $|Q(x)-L|<\varepsilon/4$. Mamy dokładną tożsamość
\[
 \frac{f(x)}{g(x)}-L
 =(Q(x)-L)\left(1-\frac{g(b)}{g(x)}\right)
 +\frac{f(b)-Lg(b)}{g(x)}.
\]
Ponieważ $b$ jest ustalone i $|g(x)|\to+\infty$, możemy ograniczyć $x$ do takiego otoczenia końca $a$, że $|g(b)/g(x)|<1/2$ oraz $|(f(b)-Lg(b))/g(x)|<\varepsilon/2$. Wtedy moduł lewej strony jest mniejszy niż $3\varepsilon/8+\varepsilon/2<\varepsilon$. Dostajemy żądaną granicę.

Jeżeli $L=+\infty$, dla $K>0$ wybieramy przedział przylegający do $a$, gdzie $g'\ne0$ i $f'/g'>2(K+1)$. Ustalamy w nim $b$. Dla $x$ między $b$ i $a$ mamy $Q(x)>2(K+1)$. Z tożsamości
\[
 \frac{f(x)}{g(x)}
 =Q(x)\left(1-\frac{g(b)}{g(x)}\right)+\frac{f(b)}{g(x)}
\]
i warunków $|g(b)/g(x)|<1/2$, $|f(b)/g(x)|<1$, które zachodzą dostatecznie blisko $a$, wynika $f(x)/g(x)>K+1-1=K$. Dla $L=-\infty$ stosujemy przypadek dodatniej nieskończoności do $-f$ i $g$. To kończy dowód.
\end{proof}
Klasyczna postać $\infty/\infty$ jest szczególnym przypadkiem tego twierdzenia. Nie wymaga ono dodatniości mianownika ani jego pochodnej: w dowodzie czynnik $1-g(b)/g(x)$ jest ostatecznie dodatni, niezależnie od znaku $g(x)$.

\section{Przykłady i przekształcanie postaci nieoznaczonych}
\begin{przyklad}
\label{ma1:ex:lhopital-examples}
Przy $x\to0$ mamy $1-\cos x\to0$ i $x^2\to0$. Na obu stronach zera pochodna mianownika $2x$ jest niezerowa. Iloraz pochodnych to $\sin x/(2x)$, który dąży do $1/2$, ponieważ $\sin$ jest różniczkowalna w zerze i $\sin'(0)=1$. Reguła \ref{ma1:thm:lhopital-zero} daje
\[
 \lim_{x\to0}\frac{1-\cos x}{x^2}=\frac12.
\]
Alternatywą jest rozwinięcie Taylora z rozdziału \ref{ma1:ch:taylor}.

Dla $x\to+\infty$ mianownik $x$ dąży do nieskończoności, jego pochodna wynosi $1$, a pochodna licznika $\log x$ wynosi $1/x\to0$. Z twierdzenia \ref{ma1:thm:lhopital-infinity} wynika $\log x/x\to0$, zgodnie z ogólniejszym twierdzeniem \ref{ma1:thm:asymptotic-growth}.

Jeżeli $f$ i $g$ spełniają założenia reguły dla $0/0$ lub nieskończonego mianownika, ale granicę $f'/g'$ chcemy obliczyć kolejnym zastosowaniem reguły, musimy osobno sprawdzić jej założenia dla nowego ilorazu. Na przykład dla $x\to0$:
\[
 \frac{e^x-1-x}{x^2}
 \quad\longrightarrow\quad \frac{e^x-1}{2x}
 \quad\longrightarrow\quad \frac{e^x}{2}.
\]
Strzałki oznaczają przejście do ilorazu pochodnych, nie równość funkcji. Pierwsze dwa ilorazy mają postać $0/0$, a pochodne ich mianowników to odpowiednio $2x$ (niezerowe poza zerem) i $2$. Ostatni iloraz ma granicę $1/2$. Stosując regułę najpierw do drugiego, potem do pierwszego ilorazu, otrzymujemy granicę $1/2$.
\end{przyklad}

\begin{przyklad}
\label{ma1:ex:lhopital-transformations}
Postać $0\cdot\infty$ można niekiedy przekształcić w iloraz. Dla $x>0$ i $x\to0^+$ mamy $x\log x=(\log x)/(1/x)$. Mianownik dąży do $+\infty$, jego pochodna $-1/x^2$ jest niezerowa, a iloraz pochodnych $(1/x)/(-1/x^2)=-x$ dąży do zera. Stąd $x\log x\to0$.

Dla potęg o dodatniej podstawie korzystamy z logarytmu: $x^x=\exp(x\log x)$ dla $x>0$, więc przy $x\to0^+$ z ciągłości funkcji wykładniczej wynika $x^x\to1$. Podobnie przy postaciach $1^\infty$ i $\infty^0$ najpierw badamy iloczyn wykładnika i logarytmu podstawy; dopiero po ustaleniu jego granicy przechodzimy do funkcji wykładniczej.

Dla postaci $\infty-\infty$ często można użyć wspólnego mianownika lub sprzężenia. Na przykład dla $x>0$,
\[
 \sqrt{x^2+x}-x=\frac{x}{\sqrt{x^2+x}+x}
 =\frac{1}{\sqrt{1+1/x}+1}\longrightarrow\frac12
 \qquad(x\to+\infty).
\]
Tutaj przekształcenie algebraiczne wystarcza i reguła de l'Hospitala nie jest potrzebna.
\end{przyklad}

\section{Granice stosowalności}
\begin{przyklad}
\label{ma1:ex:lhopital-counterexamples}
Reguła jest twierdzeniem o granicach, nie regułą różniczkowania ilorazu. Bez odpowiedniej postaci granicznej nie wolno jej stosować. Dla $x\to0$ iloraz $(1+x)/(2+x)$ dąży do $1/2$, podczas gdy iloraz pochodnych jest stale równy $1$. Mianownik dąży do $2$, a licznik do $1$, więc nie zachodzą założenia żadnej z obu wersji.

Istnienie granicy ilorazu pochodnych jest warunkiem wystarczającym w tych twierdzeniach, lecz nie koniecznym dla istnienia granicy wyjściowej. W przykładzie \ref{ma1:ex:asymptotic-cancellation} funkcje $F(x)=x+x^2\sin(1/x)$ i $G(x)=x$ na $(0,\infty)$ dążą do zera przy $x\to0^+$. Mamy $G'=1$ oraz $F/G\to1$, ale $F'/G'$ nie ma granicy. Jest to przykład właściwej postaci $0/0$, w której nie można uzyskać wyniku z reguły.

Niezerowość $g'$ w odpowiednim otoczeniu gwarantuje, że wersja ilorazowa twierdzenia Cauchy'ego jest dostępna na każdym używanym odcinku. Nie wolno pomijać tego założenia w przedstawionych dowodach ani uznawać ilorazu $f'/g'$ za określony w zerach mianownika (por. przykład \ref{ma1:ex:cauchy-degenerate}).
\end{przyklad}

Praktyczna kolejność postępowania: ustalamy granice licznika i mianownika, wybieramy odpowiedni kierunek, sprawdzamy różniczkowalność i niezerowość pochodnej mianownika, a następnie badamy granicę ilorazu pochodnych. Jeżeli potrzebne są dalsze przekształcenia lub kolejne zastosowania reguły, kontrolę przeprowadzamy ponownie. Gdy dostępne są prostsze oszacowania lub rozwinięcia asymptotyczne, warto je wykorzystać zamiast mechanicznego różniczkowania.
